\section{Introdução}

As membranas biológicas são fronteiras ativas e altamente complexas que delimitam os espaços celulares e regulam o trânsito de íons, nutrientes e sinais bioquímicos. A integridade estrutural e a funcionalidade destas interfaces dependem intimamente da sua composição lipídica e do ambiente termodinâmico em que estão inseridas. No modelo do mosaico fluido, os fosfolipídios organizam-se espontaneamente em bicamadas, cujas propriedades mecânicas e transições de fase são ditadas pelo comprimento das cadeias acílicas, pelo grau de insaturação e, fundamentalmente, pelas interações de hidratação e força iónica na interface polar.

Dentre os componentes moduladores da membrana celular, o colesterol desempenha um papel físico-químico central. A sua arquitetura molecular rígida atua de forma decisiva no empacotamento das cadeias de fosfolipídios, como a dipalmitoilfosfatidilcolina (DPPC) e a diestearoilfosfatidilcolina (DSPC). A intercalação do colesterol induz a formação da fase líquida-ordenada ($L_o$), promovendo o chamado efeito condensador, caracterizado pela redução drástica da Área por Lipídio (APL) e pelo aumento da espessura da membrana. Esta modulação é vital para a formação de microdomínios funcionais e para a regulação de processos biofísicos essenciais, incluindo a ancoragem de proteínas e a fusão de membranas.

A investigação detalhada da termodinâmica desses sistemas em escala nanométrica tem sido impulsionada pela evolução da Dinâmica Molecular (DM). Embora as simulações em nível atomístico forneçam resoluções químicas precisas, a observação de fenômenos cooperativos de grande escala — como as transições de fase lipídica e a retenção em estados metaestáveis (líquidos super-resfriados) — é severamente limitada pelos altos custos computacionais e pelas curtas escalas de tempo acessíveis. Para transpor estas barreiras cinéticas, a modelagem \textit{coarse-grained} (CG), em particular o recente campo de força Martini 3, surge como uma abordagem indispensável, permitindo o escrutínio do balanço energético global em escalas de tempo de centenas de nanossegundos.

Neste contexto, o presente trabalho emprega a simulação computacional avançada para isolar e quantificar os eixos termodinâmicos (forças de Coulomb e de Lennard-Jones) que governam o comportamento de bicamadas lipídicas puras e em misturas com colesterol, elucidando as bases energéticas da homeostase e da estabilidade estrutural das membranas celulares.

\section{OBJETIVOS}
\subsection{Objetivo Geral}
Analisar os mecanismos estruturais e termodinâmicos que permitem  a estabilidade, a transição de fase e a modulação mecânica de membranas lipídicas (DPPC e DSPC), na presença e ausência de colesterol, utilizando simulações de Dinâmica Molecular com o modelo \textit{coarse-grained} Martini 3.

\subsection{Objetivos Específicos}
\begin{itemize}
     \item Investigar a dependência do grau de hidratação na estabilidade estrutural das bicamadas puras e identificar os fatores termodinâmicosresponsáveis pela metaestabilidade observada no DSPC;
    \item Avaliar o impacto das interações eletrostáticas, em especial a coordenação do cloreto de sódio ($NaCl$), na densidade e estruturação da interface polar da membrana;
    \item Analisar o efeito condensador induzido por concentrações crescentes de colesterol nas propriedades físicas da membrana, relacionando a variação da Área por Lipídio com o balanço entálpico global da caixa de simulação.
     \item Quantificar as contribuições isoladas das energias de Coulomb e de Lennard-Jones frente às variações de temperatura e composição lipíd
   ica;
\end{itemize}
